\documentclass[11pt]{article}
\usepackage{icmlko}
\begin{document}

\title{손으로 고른 문턱과 원리적인 축이 서로를 대체하지 않는다 --- 그리고 $\tau$ 는 여덟이다}
\author{월드모델 실험실}
\date{노트 517 · 2026년 8월 4일}
\maketitle

\begin{abstract}
노트 516 이 \texttt{TRAINR}(연도 반감기 $8$ 년)을 채택하면서 물음을 남겼다
--- \emph{연도 가중이 만화(제일 옛날인 도메인)를 이미 크게 줄이는데
\texttt{TRAINCAP}\{만화$:0.5$\} 이 아직 필요한가.} 필요 없다면 손으로 고른
도메인 문턱을 빼고 원리적인 축 하나로 갈 수 있다.

\textbf{반증 --- 안 겹친다.}

\begin{center}
\begin{tabular}{lrrrrr}
\hline
현 챔피언 대비 & 판 & 날짜 통제 & 무리 안 & KR 만화 & 비게임 앱 \\
\hline
\texttt{TRAINCAP} 끔 & $\mathbf{-0.0020}$ & $\mathbf{-0.0015}$ & $\mathbf{-0.0021}$ & $-0.0023$ & $-0.0035$ \\
$\tau{=}12$ & $-0.0001$ & $+0.0002$ & $-0.0008$ & $-0.0006$ & $-0.0018$ \\
$\tau{=}16$ & $\mathbf{-0.0011}$ & $-0.0005$ & $\mathbf{-0.0017}$ & $\mathbf{-0.0042}$ & $\mathbf{-0.0041}$ \\
\hline
\end{tabular}
\end{center}

\texttt{TRAINCAP} 을 끄면 \textbf{집안 자 셋이 전부 $0$ 밖 음수}다.
$\tau$ 는 $12$ 가 동점이고 $16$ 은 확실히 나쁘다 --- \textbf{$8$ 이 봉우리
근처다.}
\end{abstract}

\begin{figure}[t]
\centering
\includegraphics[width=\textwidth]{figs/f1.pdf}
\caption{\textbf{가} --- 세 팔. \textbf{나} --- $\tau$ 훑기.
\textbf{다} --- 셋이 사는 것.}
\end{figure}

\section{셋이 서로 다른 것을 산다}

이 사이클에서 채택한 학습 쪽 손잡이 셋을 하나씩 빼 보면 각자 다른 자를
떨어뜨린다.

\begin{center}
\begin{tabular}{lll}
\hline
손잡이 & 빼면 제일 아픈 곳 & 크기 \\
\hline
\texttt{TRAINCAP}(만화 절반) & 집안 셋이 고르게 & 판 $-0.0020$ · 무리 안 $-0.0021$ \\
\texttt{TRAINW}(유보 무게) & \textbf{KR 만화} & $-0.0139$ (노트 516) \\
\texttt{TRAINR}(연도 $\tau{=}8$) & \textbf{판 $\cdot$ 날짜 통제} & 판 $-0.0025$ (노트 516) \\
\hline
\end{tabular}
\end{center}

\textbf{셋이 다 ``옛 행을 덜 본다''인데 떨어뜨리는 자가 다르다.} 겹치는
일을 한다면 하나를 빼도 다른 둘이 메워야 하는데 안 메운다.

읽으면 --- \textbf{옛 행을 줄이는 방법마다 ``어느 옛 행''이 다르다.}
\texttt{TRAINCAP} 은 만화의 \emph{무작위} 절반, \texttt{TRAINW} 는
\emph{도메인 통째로} 비례해서, \texttt{TRAINR} 은 \emph{연도순으로}.
세 방식이 남기는 표본이 서로 달라서 각자 다른 자를 산다.

\begin{quote}
\textbf{조항} --- \textbf{두 손잡이가 같은 방향으로 움직인다고 겹친다고
보지 않는다.} 겹침은 \emph{하나를 뺐을 때 다른 것이 메우는가}로만 잰다.
\end{quote}

\section{$\tau$ 는 여덟}

\begin{center}
\begin{tabular}{lrr}
\hline
$\tau$ & 판 & KR 만화 \\
\hline
없음 & $-0.0025$ & $-0.0139$ \\
$\mathbf{8}$ & --- & --- \\
$12$ & $-0.0001$ & $-0.0006$ \\
$16$ & $-0.0011$ & $-0.0042$ \\
\hline
\end{tabular}
\end{center}

$\tau{=}4$ 는 노트 516 에서 $\tau{=}8$ 보다 나빴다($\tau{=}8$ 판
$+0.0025$ 대 $\tau{=}4$ $+0.0021$, KR $-0.0018$ 대 $-0.0009$ --- 둘 다
\texttt{TRAINW} 없이 잰 값이라 직접 비교는 아니다). \textbf{$8 \sim 12$
사이에 평평한 봉우리가 있고 그 밖은 내려간다.}

\textbf{정직하게} --- $\tau{=}12$ 는 다섯 자가 다 $0$ 을 물어 ``$8$ 보다
나쁘다''고 못 말한다. \textbf{동점이다.} $8$ 을 유지하는 이유는 먼저
채택했기 때문이지 더 낫다고 보였기 때문이 아니다(노트 432 의 되돌림 조항과
같은 정신).

\section*{이 사이클의 챔피언 --- 확정}

\begin{center}
\begin{tabular}{ll}
\hline
\texttt{DOMDROP} & \texttt{\{"애니": ("trend\_",)\}} \\
\texttt{TRAINCAP} & \texttt{\{"만화": 0.5\}} \\
\texttt{TRAINW} & 유보 무게 · clip $(0.2, 1.0)$ \\
\texttt{TRAINR} & $8.0$ \\
\hline
판 & $0.4812 \to \mathbf{0.4932}$ \\
\hline
\end{tabular}
\end{center}

\section*{남은 것}

\begin{itemize}
\item \textbf{\texttt{TRAINCAP} 의 만화 $0.5$ 는 손으로 골랐다} --- $0.25
\cdot 0.75$ 는 아직 안 재 봤다. 이제 \texttt{TRAINR} 위에서 재야 한다.
\item \textbf{다른 도메인에도 \texttt{TRAINCAP}} --- 세계애니는
\texttt{TRAINW} 가 $0.61$ 로 이미 줄이는데 그 위에 문턱을 더 걸면?
\item \textbf{노트 446 거부권을 다시 읽기}(노트 516) --- 학습 쪽 후보가
그 거부권에 걸려 죽은 적이 있나.
\end{itemize}

\end{document}
